\section{Providence, Will, Destiny}

\begin{quotex}
For the earnest expectation of the creation waits for the manifestation of the sons of God. 
\flright{\textsc{Romans 8:19}}

The wise man will rule the stars. 
\flright{\textit{Hermetic maxim}}

During my whole life I have not succeeded in finding a more lucid formula and a more effective general key for understanding the evolution and history of mankind than that given by \textbf{Fabre d'Olivet}.
\flright{\textsc{Valentin Tomberg}}

Fabre d'Olivet was the first, in recent time, to sustain the remote Nordic-Arctic, hyperborean origin. For him that thesis however has less the character of a scientific hypothesis than that of the exposition of a traditional teaching, still preserved in tightly closed circles with which he was in contact. 
\flright{\textsc{Julius Evola}}
\end{quotex}

\textbf{Rene Guenon} mentioned Fabre d'Olivet many times throughout his works and dedicates a chapter on Fabre's teaching on Providence, Will, and Destiny in The Great Triad. Providence is the Will of God, or \emph{natura naturans}, nature naturing, while Destiny is the obscure will of God as \emph{natura naturata}, or nature natured.

In the Philosophical History of the Human Race\footnote{\url{https://wwww.gornahoor.net/library/Hermeneutic\_interpretation\_of\_the\_origin.pdf}} Fabre follows the Traditional understanding of man as having intellectual, psychic, and instinctive centers or spirit, soul, and body. The will, then, represents the inward and central element which unites and embraces them. That represents the microcosm while Providence, Will, and Destiny represent the macrocosm. Guenon explains that the harmony between Will and providence constitutes the Good while evil is born from their opposition

\begin{quotex}
Man either perfects himself or becomes depraved according to whether his tendency is to merges into the universal Unity or to distinguish himself from it. Either he allies his will with Providence and follows the path of freedom or he allies his will with destiny and follows the path of necessity. The divine man lives a life primarily of the intellect, which is governed by providential law.
\end{quotex}


Paradoxically, then, the ``divine man'' is free because he follows the cosmic law, while the natural man believes he is free while he is under the necessity of Destiny.

\textbf{Alfred North Whitehead} in \emph{Process and Reality} illustrates this triad, if we overlook some errors. At every moment, the ``actual occasion'' is faced with two forces. The accumulated past is a weight that tends to limit possibilities. On the other side, God is providing a creative option that works to break the weight of the past. The occasion needs to integrate all that. Unfortunately, Whitehead confuses ``God'' with Providence which leads him to postulate a finite or limited god. Fabre rejects this, since God can never be subject to such conditions:

\begin{quotex}
The three powers constitute the universal ternary. Nothing escapes their action: all in the universe is subject to them; all except God himself who, enveloping them in His unfathomable unity, forms with them the sacred tetrad of the ancients, that immense quaternary, which is all in all and outside of which there is nothing.
\end{quotex}


The following schema illustrates these powers. The three centers are driven by instinct, passion, and inspiration. The Will is the heart of man when all his centers are in harmony. That is, for example, when his passions and instincts are ruled by the intellect. And conversely, when the intellect supports the true passions and instincts.

\begin{table}[h]\small\centering
\begin{tabular}{p{.3\textwidth}p{.3\textwidth}p{.3\textwidth}}
\toprule
~ & \textbf{Will}: Determination & \textbf{Providence} \\
~ & \textbf{Intellectual}: Inspiration & ~ \\
~ & \textbf{Psychic}: Passion & ~ \\
Destiny (psychic) & \textbf{Instinctive}: Instinct & ~ \\
\textbf{Destiny} (physical) & \textbf{Physical}: Necessity & ~ \\\bottomrule
\end{tabular}
\end{table}

Unfortunately, such as we are, we are not in such harmony. Tomberg relates the triad to the problem of births. Destiny is our heredity, Will is our birth, and Providence is related to salvation. Hence, that is the first task of providence.

\textbf{Pope John Paul II} in the \emph{Theology of the Body} explains that Adam lost what did not belong to human nature in the strict sense: viz., integrity, holiness, innocence, justice. In other words, these are the qualities of Fabre's ``divine man''. Man is left with his lower intellect or ``\emph{ratio}'', and therefore a weakened will. The threshold of man's earthly history is crossed with the knowledge of good and evil. Before that, there is only prehistory that is opaque to us.

\textbf{Boris Mouravieff} is in agreement with this. Adam's fall led to the loss of higher human nature, leaving him with just the lower. Specifically, he is subject to the ``general law'', like the animals, motivated primarily by hunger, sex, and fear. He is left with the ``ratio'', or lower intellect. He explains that the knowledge of good and evil can only be dualistic knowledge, the knowledge of phenomena, not the transcendental knowledge or gnosis of the higher intellectual center.

This is the Matrix we are caught in, the situation of original sin, as described by \textbf{Joseph Ratzinger}. There is interminable debate without resolution, and the constant battle of good and evil, like a point stuck inside a square in Flatland with no way out. However, if the point became aware of a third dimension, he would easily escape his trap.

It is worthwhile, then, to investigate the qualities that have been lost in the Matrix. For those concerned about tradition, they are all taken from documents and councils over several centuries.

\textbf{Innocence}: In the state of innocence, man does not know good and evil, but rather he knows by direction intuition. His heart is the tranquil witness of consciousness.

\textbf{Holiness}: Holiness is the irradiation of the Holy Spirit, which produces a special state of spiritualization, i.e., participation in God's inner life. If the lower centers are disordered, they will not reflect that irradiation accurately.

\textbf{Justice:} The just man is in conformance with the Logos or covenant law; he judges impartially.

\textbf{Integrity}: Integrity refers to man as a whole. JPII writes that the heart is the center of man understood as the source of will, emotion, thoughts, and affections. Or, if Fabre's system, all the centers are in harmonious relationship. Then there can be the determination of the True Will.


\flrightit{Posted on 2015-05-13 by Cologero}

\begin{center}* * *\end{center}

\begin{footnotesize}
\begin{sffamily}
\texttt{Wayne Ferguson on 2015-05-14 at 08:39 said:}

I had occasion to dig out an old copy of Kierkeggard's ``Sickness Unto Death'' yesterday and (after reading this last night) just happened to open it up to this text which I had marked when I first read it 23 years ago:

`` People do not generally regard being under a delusion as the greatest misfortune; their sensuous nature is generally predominant over their intellectuality. So when a man is supposed to be happy, he imagines that he is happy (whereas viewed in the light of the truth he is unhappy), and in this case he is generally very far from wishing to be torn away from that delusion. On the contrary, he becomes furious, he regards the man who does this as his most spiteful enemy, he considers it an insult, something near to murder, in the sense that one speaks of killing joy. What is the reason of this? The reason is that the sensuous nature and the psycho-sensuous completely dominate him; the reason is that he lives in the sensuous categories agreeable/disagreeable, and says goodbye to truth etc.; the reason is that he is too sensuous to have the courage to venture to be spirit or to endure it. However vain and conceited men may be, they have nevertheless for the most part a very lowly conception of themselves, that is to say, they have no conception of being spirit, the absolute of all that a man can be --- but vain and conceited they are . . . by way of comparison. In case one were to think of a house, consisting of cellar, ground-floor and premier étage, so tenanted, or rather so arranged, that it was planned for a distinction of rank between the dwellers on the several floors; and in case one were to make a comparison between such a house and what it is to be a man --- then unfortunately this is the sorry and ludicrous condition of the majority of men, that in their own house they prefer to live in the cellar. The soulish-bodily synthesis in every man is planned with a view to being spirit, such is the building; but the man prefers to dwell in the cellar, that is, in the determinants of sensuousness. And not only does he prefer to dwell in the cellar; no, he loves that to such a degree that he becomes furious if anyone would propose to him to occupy the bel étage which stands empty at his disposition --- for in fact he is dwelling in his own house.''

\url{http://www.naturalthinker.net/trl/texts/Kierkegaard,Soren/TheSicknessUntoDeath.pdf}


A quote from Peter Kingley also comes to mind

``We are human beings, endowed with an incredible dignity; but there's nothing more undignified than forgetting our greatness and clutching at straws'' (Dark Places of Wisdom 4).

\url{https://www.peterkingsley.org/cw3/Admin/images/SpiritualTradition.pdf}


\hfill

\texttt{Iron Man on 2015-05-14 at 11:27 said:}

``Man either perfects himself or becomes depraved''

Like when man is given a web site he immediately turns into a mini tyrant controlling his followers' minds and stamping out all dissent. Afraid of every last opinion, he becomes surrounded only with echoes of himself.

``Fabre d'Olivet was the first, in recent time, to sustain the remote Nordic-Arctic, hyperborean origin. For him that thesis however has less the character of a scientific hypothesis than that of the exposition of a traditional teaching, still preserved in tightly closed circles with which he was in contact.'' Julius Evola

Finally we are revealing that many of Guenon's and Evola's sources are highly suspect, as are a good deal of his theories based on nothing more than the fairytales of occultists.

\hfill

\texttt{Cologero on 2015-05-14 at 12:52 said:}

Iron Man, you sure are persistent. Blake said that if a fool persists in his folly he may become wise. Unfortunately, I've heard that there have been many exceptions.

This web site is my personal domain and you therefore have the status of a guest. If a gentleman visits a friend, he will accept the hospitality and any fare offered with good grace. Unfortunately, you are crude, vulgar, and of low breeding, unfamiliar with the protocols of hospitality. Hospitality does not mean ``control''. Rather is means friends will gather together to share thoughts and ideas. Those who are of a different mind will move on.

I'm afraid that most of us, including readers, are awash in a sea of dissent all day long ... from family, friends, the media, etc. Why would you object that like-minded people get together, since we find so little support for our thoughts in the world?

I'm flattered in an odd way that you think I can control minds. So I will let you in on a little secret. Those who can be controlled barely have a mind to speak of, at least not in the fullness we have described in these pages. And those with actual minds cannot be controlled.

While I'm intrigued that you have spent so much time reading Gornahoor over the years, I'd prefer that you remain silent, at least until you learn both the meaning of hospitality and how to write an intelligent dissent. Specificity is what I'm looking for.

\hfill

\texttt{Mark Citadel on 2015-05-17 at 11:57 said:}

Excellent and so much needed in an age where too many Christians, including myself up until recently, are unsure of our constituent parts that transcend biology, the difference between soul and spirit. Presumably the lost higher characteristics you outline here will be returned upon the post-apocalyptic resurrection?

@Ironman -- What do you find objectionable about D'Olivet? Admittedly I didn't know much about him prior to now (thanks for the heads up on this guy by the way, Cologero), but he doesn't seem to be ``highly suspect''. Yes, he may have made some mistakes during his pre-Rosetta studies of hieroglyphs which angered many people including the Pope, but that doesn't discredit his work on the pre-Fall nature of man. In fact I'd say he's one of the more thoughtful informants of Guenon and Evola's work.

He even apparently attacked Lord Byron for an anti-Christian play that he felt would undermine the Faith. Byron's response: ``I'm only a poet; I don't know anything about these philosophical concerns of yours!''

I digress, in a sane society, such a stupid person wouldn't be allowed to produce a pamphlet, let alone a play.

\end{sffamily}
\end{footnotesize}

